\documentclass[12pt,letterpaper]{article}

\usepackage[spanish]{babel}
\usepackage[utf8]{inputenc}
\usepackage[T1]{fontenc}
\usepackage{graphicx}
\usepackage{fancyhdr}
\usepackage{geometry}
\usepackage{amsmath, amssymb}

\geometry{margin=2.5cm}
\setlength{\headheight}{52pt}
\pagestyle{fancy}
\fancyhf{}

\fancyhead[L]{
    \includegraphics[height=1.5cm]{ur_image.png}
}

\fancyhead[R]{
\begin{tabular}{r}
\textbf{Monitoria Álgebra Lineal} \\
Gabriel Fernando Salcedo Zamora \\
\end{tabular}
}

\renewcommand{\headrulewidth}{0.4pt}

\begin{document}

\section*{Taller de Matrices}

\begin{enumerate}

    %producto punto
    
    \item En los siguientes problemas calcule el producto punto de los vectores. Indique además en cuáles casos los vectores son ortogonales.

    \begin{enumerate}
        \item $(2, -1) \cdot (4, 5)$
        \item $(3,-2,5) \cdot (-1,-2,6)$
        \item $(\pi, \frac{\pi^2}{3}, 3) \cdot (\pi^2, -9\pi, \pi^3)$
        \item $(x, y, z) \cdot (y, z, x)$
    \end{enumerate}

    %vectores unitarios

    \item Dado $\mathbf{u}=(2,4,2,-1)$, obtenga un vector $\mathbf{w}$ que sea paralelo al vector $\mathbf{u}$ tal que $\|w\|=10$
    
    %operaciones de matrices
    
    \item Resuelva los siguientes operaciones teniendo en cuenta que: 
    \[
    A = \begin{pmatrix} 1 & 0 & 6\\4 & 1 & 1 \\ -2 & 9 & 2 \end{pmatrix},
    \quad 
    B = \begin{pmatrix} -2 & 5 & -9\\3 & -4 & 1 \\ -1 & 0 & -6 \end{pmatrix},
    \quad
    C = \begin{pmatrix} 7 & 4 & 2 \\ -5 & -2 & 3 \\ 1 & 5 & 7     \end{pmatrix}
    \]
    \begin{enumerate}
        \item $A-2B$
        \item $A+B+C$
        \item $3B-2A$
        \item $3A+2B-4C$
    \end{enumerate}
    
    %multiplicación de matrices
    
    \item  Calcule el producto de las siguientes matrices o explique por qué no es posible.
    \begin{enumerate}
        \item $\begin{pmatrix} 1 & -1 \\ 1 & 1 \end{pmatrix}\begin{pmatrix} -1 & 0 \\ 2 & 3   \end{pmatrix}$
        \item $\begin{pmatrix} 2 & -1\end{pmatrix}\begin{pmatrix} 1 & 0 \\ 0 & 1   \end{pmatrix}$
        \item $\begin{pmatrix} 1 & 4 & 0 & 2 \end{pmatrix}\begin{pmatrix} 3 & -6 \\ 2 & 4 \\ 1 & 0 \\ -2 & 3  \end{pmatrix}$
        \item $\begin{pmatrix} a & b & c \\ d & e & f \\ g & h & j \end{pmatrix}\begin{pmatrix} 1 & 0 & 0 \\ 0 & 1 & 0 \\ 0 & 0 & 1 \end{pmatrix}$, donde $a,b,c,d,e,f,g,h,j$, son números reales.
    \end{enumerate}
    
\end{enumerate}

%\section*{Solución}
\newpage

\section*{Solución}

\begin{enumerate}

%-------------------------------------------------
% Producto punto
%-------------------------------------------------

\item Producto punto y ortogonalidad.

\begin{enumerate}
    \item 
    \[
    (2,-1)\cdot(4,5)=2(4)+(-1)(5)=8-5=3
    \]
    Como el producto punto es distinto de cero, los vectores no son ortogonales.

    \item 
    \[
    (3,-2,5)\cdot(-1,-2,6)=3(-1)+(-2)(-2)+5(6)
    \]
    \[
    =-3+4+30=31
    \]
    No son ortogonales.

    \item 
    \[
    (\pi,\tfrac{\pi^2}{3},3)\cdot(\pi^2,-9\pi,\pi^3)
    \]
    \[
    =\pi(\pi^2)+\frac{\pi^2}{3}(-9\pi)+3(\pi^3)
    \]
    \[
    =\pi^3-3\pi^3+3\pi^3=\pi^3
    \]
    No son ortogonales.

    \item 
    \[
    (x,y,z)\cdot(y,z,x)=xy+yz+zx
    \]
    Los vectores son ortogonales cuando
    \[
    xy+yz+zx=0
    \]
\end{enumerate}

%-------------------------------------------------
% Vector unitario
%-------------------------------------------------

\item Vector paralelo con norma 10.

Dado el vector $\mathbf{u}=(2,4,2,-1)$, primero calculamos su norma:
\[
\|\mathbf{u}\|=\sqrt{2^2+4^2+2^2+(-1)^2}=\sqrt{25}=5
\]

El vector unitario asociado es:
\[
\hat{u}=\frac{1}{5}(2,4,2,-1)
\]

Entonces, un vector paralelo a $\mathbf{u}$ con norma 10 es:
\[
\mathbf{w}=10\hat{u}=(4,8,4,-2)
\]

%-------------------------------------------------
% Operaciones con matrices
%-------------------------------------------------
\newpage

\item Operaciones con matrices.

\begin{enumerate}
    \item 
    \[
    A-2B=
    \begin{pmatrix}
    1 & 0 & 6\\
    4 & 1 & 1\\
    -2 & 9 & 2
    \end{pmatrix}
    -
    2\begin{pmatrix}
    -2 & 5 & -9\\
    3 & -4 & 1\\
    -1 & 0 & -6
    \end{pmatrix}
    \]
    \[
    =
    \begin{pmatrix}
    5 & -10 & 24\\
    -2 & 9 & -1\\
    0 & 9 & 14
    \end{pmatrix}
    \]

    \item 
    \[
    A+B+C=
    \begin{pmatrix}
    6 & 9 & -1\\
    2 & -5 & 5\\
    -2 & 14 & 3
    \end{pmatrix}
    \]

    \item 
    \[
    3B-2A=
    \begin{pmatrix}
    -8 & 15 & -39\\
    1 & -14 & 1\\
    1 & -18 & -22
    \end{pmatrix}
    \]

    \item 
    \[
3A+2B-4C=
\begin{pmatrix}
-29 & -6 & -8\\
38 & 3 & -7\\
-12 & 7 & -34
\end{pmatrix}
\]

\end{enumerate}

%-------------------------------------------------
% Multiplicación de matrices
%-------------------------------------------------

\item Multiplicación de matrices.

\begin{enumerate}
    \item 
    \[
    \begin{pmatrix}
    1 & -1\\
    1 & 1
    \end{pmatrix}
    \begin{pmatrix}
    -1 & 0\\
    2 & 3
    \end{pmatrix}
    =
    \begin{pmatrix}
    -3 & -3\\
    1 & 3
    \end{pmatrix}
    \]

    \item 
    \[
    \begin{pmatrix}
    2 & -1
    \end{pmatrix}
    \begin{pmatrix}
    1 & 0\\
    0 & 1
    \end{pmatrix}
    =
    \begin{pmatrix}
    2 & -1
    \end{pmatrix}
    \]

    \item 
    \[
    \begin{pmatrix}
    1 & 4 & 0 & 2
    \end{pmatrix}
    \begin{pmatrix}
    3 & -6\\
    2 & 4\\
    1 & 0\\
    -2 & 3
    \end{pmatrix}
    =
    \begin{pmatrix}
    7 & 16
    \end{pmatrix}
    \]

    \item 
    \[
    \begin{pmatrix}
    a & b & c\\
    d & e & f\\
    g & h & j
    \end{pmatrix}
    \begin{pmatrix}
    1 & 0 & 0\\
    0 & 1 & 0\\
    0 & 0 & 1
    \end{pmatrix}
    =
    \begin{pmatrix}
    a & b & c\\
    d & e & f\\
    g & h & j
    \end{pmatrix}
    \]
    El producto es la misma matriz, ya que se multiplica por la matriz identidad.
\end{enumerate}

\end{enumerate}

\end{document}
